% =====================================================================
%  related-work.tex
%  Drop into memoe.tex with % =====================================================================
%  related-work.tex
%  Drop into memoe.tex with % =====================================================================
%  related-work.tex
%  Drop into memoe.tex with % =====================================================================
%  related-work.tex
%  Drop into memoe.tex with \input{related-work.tex}
%  Suggested placement: after the serving-loop section, before
%  "Two smaller results". By that point every measurement this section
%  refers to has already been presented.
% =====================================================================

\section{How this compares to existing expert offloading}

Expert offloading is a well-populated area. Every system below starts from the
same premise we do, that expert weights dominate the footprint and should not
occupy HBM, and then diverges on what to do about it. The divergence is almost
entirely about the batch size each system was designed for, which is the axis our
closed form describes. This section places MEMoE-RT on that axis and states what
its architecture buys as the deployment grows.

\subsection{The design space, ordered by batch size}

\textbf{Predict and cache, at batch one.} The largest group treats expert
selection as predictable and expert reuse as temporally local.
Mixtral-offloading keeps an LRU cache of individual experts on the GPU and
speculatively prefetches one to two layers ahead, on the premise that consecutive
tokens activate overlapping experts (Eliseev and Mazur, arXiv:2312.17238).
MoE-Infinity adds request-level activation tracing, ProMoE trains a predictor,
HOBBIT varies expert precision by importance and compares LRU against
sequence-level LFU and a layer-distance policy, and AdapMoE skips experts judged
unimportant. All of them are built and evaluated at or near single-sequence
decoding on consumer hardware.

\textbf{Move the activation instead of the weight.} Fiddler observes that at
small batch an activation is far smaller than an expert's weights, and therefore
copies activations to the host and executes the expert layer on the CPU rather
than moving weights to the GPU (Kamahori et al., ICLR 2025, arXiv:2402.07033).
Initial placement is popularity-ranked from calibration data. The paper is
explicit that the advantage holds when the batch is sufficiently small, which
makes its design point and ours two sides of the same boundary.

\textbf{Pipeline for throughput.} ZeRO-Inference offloads layer weights to host
or NVMe with prefetch and targets throughput at large batch rather than
single-request latency; its own reported finding is that in that regime full
offload beats partial offload, because the larger batch it enables more than pays
for the transfer (DeepSpeed, 2022). FlexGen schedules tensor placement across
GPU, CPU and disk. MoE-Lightning adds a CPU-GPU-I/O pipeline with paged weights
and a Hierarchical Roofline Model used to select a policy, reporting up to
10.3$\times$ over prior offloading systems for Mixtral-8x7B on a single T4 (Cao
et al., ASPLOS 2025, arXiv:2411.11217). MoE-Gen batches per module.

\textbf{Static placement, in production.} llama.cpp's \texttt{--cpu-moe},
\texttt{--n-cpu-moe} and \texttt{--override-tensor} pin whole routed-expert
tensors to host memory by layer. There is no ranking and no prediction. It is the
simplest arrangement that works and it is what most practitioners actually run.

\subsection{Why the first two groups do not carry to serving scale}

The systems in the first two groups are not underdesigned. They are correctly
designed for $B < B^{*}$, and three of our measurements explain why their
mechanisms stop paying above it.

\emph{Prediction has nothing left to predict.} Coverage saturates at 97.3\% to
100\% by a batch of 1{,}024 tokens (Section~4.2). A speculative prefetcher that
guesses which experts will fire is, at serving batch sizes, guessing that all of
them will, and it is right. The prediction machinery becomes a fixed cost against
a decision that has become trivial.

\emph{Ranking has nothing left to rank.} Fetch hit rate equals resident fraction
exactly at every routing skew we tested at batch 512 and above
(Table~\ref{tab:skewsweep}), which collapses LRU, sampled LFU and static
placement to the same number. An LRU cache above $B^{*}$ is a fixed partition
that pays eviction bookkeeping for nothing. Worse, a placement profiled on one
domain scores 0.383 against 0.500 for random when it meets another
(Section~4.4), so calibration-derived ranking is not merely inert at scale but
can be negative.

\emph{Activation size overtakes weight size.} Fiddler's premise is that
$B \times d_{\text{model}}$ is small against $3 \, d_{\text{model}}
d_{\text{ff}}$. That ratio is linear in $B$, so the premise inverts at a batch
size determined by the same quantities that appear in Equation~\ref{eq:bstar}.
Above the crossover, moving the weights once per layer and computing on the GPU
is the cheaper arrangement, which is what MEMoE-RT does.

\emph{Per-expert granularity fails a concurrency requirement.} Systems that fetch
experts individually issue many small transfers. Section~12.3 gives the standing
requirement: saturating a tier of bandwidth $BW$ at latency $L$ with requests of
size $S$ needs $BW \cdot L / S$ requests in flight, roughly 160 for a 76.8~GB/s
expander at 64-byte granularity. Our transfer-size measurement gives the same
result from the other direction, with time spent waiting rather than moving rising
from 0.16\% for a whole 12~MB expert to 83.18\% at 4~KB. A fetch engine built
around individual experts is committed to a concurrency budget it must then work
to meet.

\subsection{What our architecture does instead}

MEMoE-RT is built from those three findings rather than around them, and each
design decision is a deletion.

Because coverage saturates, there is no predictor. Because ranking is inert, there
is no popularity ordering anywhere in the implementation; the resident set is the
low expert indices. Because latency is negligible at expert granularity and
dominant below it, there is no per-expert fetch path: every non-resident expert
of layer $L+d$ moves as one contiguous transfer, issued $d$ layers ahead on a
dedicated copy stream, into a ring of $d+1$ staging buffers where layer $L$
occupies slot $L \bmod (d+1)$.

That single-descriptor-per-layer structure is what makes the concurrency
requirement disappear rather than be satisfied: one large DMA replaces thousands
of in-flight small reads, and a single expert-sized transfer already reaches 90\%
of the PCIe plateau. It is also what made deep lookahead unnecessary. Our own
simulator, which modelled per-expert on-demand fetches, recommended depth 4 to 8;
the runtime measured depth 1 as sufficient and depth 8 as 41\% more staging memory
for slightly less throughput (Section~7.3). Batching the transfer by layer removed
the need for depth, and we report the correction rather than the original
recommendation.

The second architecture path is the part that generalises furthest. Rather than
reimplement routing per model family, it points each offloaded expert's projection
weights at a view into the staging buffer and drives the tier from pre- and
post-hooks on the MoE block. Because layer $L$ always lands in slot
$L \bmod (d{+}1)$ at a fixed offset, those views are bound once at load time and
remain correct for the process lifetime. The model's forward is untouched. This is
why the same runtime covers both a fused-expert checkpoint and DeepSeek-V2-Lite's
own modelling code, with different top-$k$, shared experts, a dense first layer
and multi-head latent attention, without architecture-specific transfer logic.

\subsection{Scaling behaviour}

The distinction that matters for a deployment is how each design behaves as the
workload grows along four axes.

\begin{table}[htbp]
\centering\small
\caption{Scaling behaviour by axis. The first two groups are built for the
left-hand end of the batch axis and their mechanisms lose value moving right; the
throughput group and MEMoE-RT gain value moving right.}
\label{tab:scaling}
\begin{tabular}{lll}
\toprule
Axis & Predict-and-cache designs & MEMoE-RT \\
\midrule
Batch size & Hit rate $\to$ residency; predictor idles & Cost per token falls as $1/B$ \\
Transfer count & Grows with experts touched & One descriptor per layer, fixed \\
Model size & Bounded by what the GPU holds & Bounded by host memory \\
GPU count & Placement re-profiled per deployment & $B^{*}$ independent of $N$ \\
\bottomrule
\end{tabular}
\end{table}

\textbf{Batch.} A fetched expert serves $Bk/E$ tokens, so expert arithmetic
intensity grows with the batch while the transfer cost is paid once. This is why
the same configuration costs 62\% of throughput at a decode-sized batch and 5.6\%
at a prefill-sized one, and why offloaded decode time per output token is flat at
$325 \pm 2$~ms across an eightfold change in concurrency while resident cost grows
linearly (Section~10.2).

\textbf{Transfer count.} Per-layer batching makes the number of transfers a
function of depth and layer count alone, not of routing, batch size or residency.
The bytes moved scale; the descriptor count does not.

\textbf{Model size.} Because placement is by capacity rather than by profile, the
resident fraction is a free parameter down to zero. DeepSeek-V2-Lite is 31.4~GB
against a 23~GB device, so no resident baseline exists at all, and MEMoE-RT runs
it in 15.13~GB at 7{,}241 tokens per second. Designs that keep a hot set sized
from a profile do not have that degree of freedom in the same way.

\textbf{GPU count.} $B^{*}$ is independent of both expert size and GPU count,
because a larger expert costs proportionally more to move and does proportionally
more work once moved, and because compute and per-node bandwidth scale together
under tensor parallelism. The operating point therefore transfers across
deployment sizes without re-derivation. Pooling is the one place this changes, and
it changes predictably: capacity saving and required batch size both scale with
node count, at one to one.

\subsection{Where the contribution sits}

Several of MEMoE-RT's mechanisms have precedent, and the comparison is more useful
for saying so. Layer-granular prefetch on a separate stream is what
ZeRO-Inference does. Static unranked expert placement is what llama.cpp does. That
full offload beats partial offload in the throughput regime is ZeRO-Inference's
result, not ours. Having an analytical model select a policy is MoE-Lightning's
HRM.

What is new is the boundary itself. Existing systems optimise a point in the
design space and report a speedup against a chosen baseline; Equation~\ref{eq:bstar}
gives the batch size at which expert offload stops costing throughput, in closed
form, independent of expert size and GPU count, and predicts the measured
crossover to within 0.5\% on OLMoE-1B-7B and to within the resolution of our
measurements on DeepSeek-V2-Lite. MoE-Lightning's HRM is the nearest relative and
selects a micro-batch policy for a system in hand; $B^{*}$ is a threshold a reader
can evaluate for hardware they do not own.

The second new element is the reason behind llama.cpp's design. Unranked static
placement is, on our measurements, the correct choice at serving scale, and we can
now say why: hit rate equals residency, so the ranking that the rest of the
literature computes buys nothing above $B^{*}$, and a ranking computed on the
wrong domain buys less than nothing. What remains is a capacity decision with a
closed form, and a scheduling problem solved by moving whole layers one step ahead
of demand.

These design points are complementary rather than competing, and they partition
cleanly by regime: predict-and-cache below $B^{*}$, contiguous layer streaming
above it. Published figures across the two groups are not directly comparable,
since retention against a resident baseline on a GPU that fits the model measures
a different quantity from speedup against another offloading system on a GPU that
does not. Running these systems on common hardware across the batch range, with
$B^{*}$ predicting where each should win, is the experiment this framework is
built to support.

%  Suggested placement: after the serving-loop section, before
%  "Two smaller results". By that point every measurement this section
%  refers to has already been presented.
% =====================================================================

\section{How this compares to existing expert offloading}

Expert offloading is a well-populated area. Every system below starts from the
same premise we do, that expert weights dominate the footprint and should not
occupy HBM, and then diverges on what to do about it. The divergence is almost
entirely about the batch size each system was designed for, which is the axis our
closed form describes. This section places MEMoE-RT on that axis and states what
its architecture buys as the deployment grows.

\subsection{The design space, ordered by batch size}

\textbf{Predict and cache, at batch one.} The largest group treats expert
selection as predictable and expert reuse as temporally local.
Mixtral-offloading keeps an LRU cache of individual experts on the GPU and
speculatively prefetches one to two layers ahead, on the premise that consecutive
tokens activate overlapping experts (Eliseev and Mazur, arXiv:2312.17238).
MoE-Infinity adds request-level activation tracing, ProMoE trains a predictor,
HOBBIT varies expert precision by importance and compares LRU against
sequence-level LFU and a layer-distance policy, and AdapMoE skips experts judged
unimportant. All of them are built and evaluated at or near single-sequence
decoding on consumer hardware.

\textbf{Move the activation instead of the weight.} Fiddler observes that at
small batch an activation is far smaller than an expert's weights, and therefore
copies activations to the host and executes the expert layer on the CPU rather
than moving weights to the GPU (Kamahori et al., ICLR 2025, arXiv:2402.07033).
Initial placement is popularity-ranked from calibration data. The paper is
explicit that the advantage holds when the batch is sufficiently small, which
makes its design point and ours two sides of the same boundary.

\textbf{Pipeline for throughput.} ZeRO-Inference offloads layer weights to host
or NVMe with prefetch and targets throughput at large batch rather than
single-request latency; its own reported finding is that in that regime full
offload beats partial offload, because the larger batch it enables more than pays
for the transfer (DeepSpeed, 2022). FlexGen schedules tensor placement across
GPU, CPU and disk. MoE-Lightning adds a CPU-GPU-I/O pipeline with paged weights
and a Hierarchical Roofline Model used to select a policy, reporting up to
10.3$\times$ over prior offloading systems for Mixtral-8x7B on a single T4 (Cao
et al., ASPLOS 2025, arXiv:2411.11217). MoE-Gen batches per module.

\textbf{Static placement, in production.} llama.cpp's \texttt{--cpu-moe},
\texttt{--n-cpu-moe} and \texttt{--override-tensor} pin whole routed-expert
tensors to host memory by layer. There is no ranking and no prediction. It is the
simplest arrangement that works and it is what most practitioners actually run.

\subsection{Why the first two groups do not carry to serving scale}

The systems in the first two groups are not underdesigned. They are correctly
designed for $B < B^{*}$, and three of our measurements explain why their
mechanisms stop paying above it.

\emph{Prediction has nothing left to predict.} Coverage saturates at 97.3\% to
100\% by a batch of 1{,}024 tokens (Section~4.2). A speculative prefetcher that
guesses which experts will fire is, at serving batch sizes, guessing that all of
them will, and it is right. The prediction machinery becomes a fixed cost against
a decision that has become trivial.

\emph{Ranking has nothing left to rank.} Fetch hit rate equals resident fraction
exactly at every routing skew we tested at batch 512 and above
(Table~\ref{tab:skewsweep}), which collapses LRU, sampled LFU and static
placement to the same number. An LRU cache above $B^{*}$ is a fixed partition
that pays eviction bookkeeping for nothing. Worse, a placement profiled on one
domain scores 0.383 against 0.500 for random when it meets another
(Section~4.4), so calibration-derived ranking is not merely inert at scale but
can be negative.

\emph{Activation size overtakes weight size.} Fiddler's premise is that
$B \times d_{\text{model}}$ is small against $3 \, d_{\text{model}}
d_{\text{ff}}$. That ratio is linear in $B$, so the premise inverts at a batch
size determined by the same quantities that appear in Equation~\ref{eq:bstar}.
Above the crossover, moving the weights once per layer and computing on the GPU
is the cheaper arrangement, which is what MEMoE-RT does.

\emph{Per-expert granularity fails a concurrency requirement.} Systems that fetch
experts individually issue many small transfers. Section~12.3 gives the standing
requirement: saturating a tier of bandwidth $BW$ at latency $L$ with requests of
size $S$ needs $BW \cdot L / S$ requests in flight, roughly 160 for a 76.8~GB/s
expander at 64-byte granularity. Our transfer-size measurement gives the same
result from the other direction, with time spent waiting rather than moving rising
from 0.16\% for a whole 12~MB expert to 83.18\% at 4~KB. A fetch engine built
around individual experts is committed to a concurrency budget it must then work
to meet.

\subsection{What our architecture does instead}

MEMoE-RT is built from those three findings rather than around them, and each
design decision is a deletion.

Because coverage saturates, there is no predictor. Because ranking is inert, there
is no popularity ordering anywhere in the implementation; the resident set is the
low expert indices. Because latency is negligible at expert granularity and
dominant below it, there is no per-expert fetch path: every non-resident expert
of layer $L+d$ moves as one contiguous transfer, issued $d$ layers ahead on a
dedicated copy stream, into a ring of $d+1$ staging buffers where layer $L$
occupies slot $L \bmod (d+1)$.

That single-descriptor-per-layer structure is what makes the concurrency
requirement disappear rather than be satisfied: one large DMA replaces thousands
of in-flight small reads, and a single expert-sized transfer already reaches 90\%
of the PCIe plateau. It is also what made deep lookahead unnecessary. Our own
simulator, which modelled per-expert on-demand fetches, recommended depth 4 to 8;
the runtime measured depth 1 as sufficient and depth 8 as 41\% more staging memory
for slightly less throughput (Section~7.3). Batching the transfer by layer removed
the need for depth, and we report the correction rather than the original
recommendation.

The second architecture path is the part that generalises furthest. Rather than
reimplement routing per model family, it points each offloaded expert's projection
weights at a view into the staging buffer and drives the tier from pre- and
post-hooks on the MoE block. Because layer $L$ always lands in slot
$L \bmod (d{+}1)$ at a fixed offset, those views are bound once at load time and
remain correct for the process lifetime. The model's forward is untouched. This is
why the same runtime covers both a fused-expert checkpoint and DeepSeek-V2-Lite's
own modelling code, with different top-$k$, shared experts, a dense first layer
and multi-head latent attention, without architecture-specific transfer logic.

\subsection{Scaling behaviour}

The distinction that matters for a deployment is how each design behaves as the
workload grows along four axes.

\begin{table}[htbp]
\centering\small
\caption{Scaling behaviour by axis. The first two groups are built for the
left-hand end of the batch axis and their mechanisms lose value moving right; the
throughput group and MEMoE-RT gain value moving right.}
\label{tab:scaling}
\begin{tabular}{lll}
\toprule
Axis & Predict-and-cache designs & MEMoE-RT \\
\midrule
Batch size & Hit rate $\to$ residency; predictor idles & Cost per token falls as $1/B$ \\
Transfer count & Grows with experts touched & One descriptor per layer, fixed \\
Model size & Bounded by what the GPU holds & Bounded by host memory \\
GPU count & Placement re-profiled per deployment & $B^{*}$ independent of $N$ \\
\bottomrule
\end{tabular}
\end{table}

\textbf{Batch.} A fetched expert serves $Bk/E$ tokens, so expert arithmetic
intensity grows with the batch while the transfer cost is paid once. This is why
the same configuration costs 62\% of throughput at a decode-sized batch and 5.6\%
at a prefill-sized one, and why offloaded decode time per output token is flat at
$325 \pm 2$~ms across an eightfold change in concurrency while resident cost grows
linearly (Section~10.2).

\textbf{Transfer count.} Per-layer batching makes the number of transfers a
function of depth and layer count alone, not of routing, batch size or residency.
The bytes moved scale; the descriptor count does not.

\textbf{Model size.} Because placement is by capacity rather than by profile, the
resident fraction is a free parameter down to zero. DeepSeek-V2-Lite is 31.4~GB
against a 23~GB device, so no resident baseline exists at all, and MEMoE-RT runs
it in 15.13~GB at 7{,}241 tokens per second. Designs that keep a hot set sized
from a profile do not have that degree of freedom in the same way.

\textbf{GPU count.} $B^{*}$ is independent of both expert size and GPU count,
because a larger expert costs proportionally more to move and does proportionally
more work once moved, and because compute and per-node bandwidth scale together
under tensor parallelism. The operating point therefore transfers across
deployment sizes without re-derivation. Pooling is the one place this changes, and
it changes predictably: capacity saving and required batch size both scale with
node count, at one to one.

\subsection{Where the contribution sits}

Several of MEMoE-RT's mechanisms have precedent, and the comparison is more useful
for saying so. Layer-granular prefetch on a separate stream is what
ZeRO-Inference does. Static unranked expert placement is what llama.cpp does. That
full offload beats partial offload in the throughput regime is ZeRO-Inference's
result, not ours. Having an analytical model select a policy is MoE-Lightning's
HRM.

What is new is the boundary itself. Existing systems optimise a point in the
design space and report a speedup against a chosen baseline; Equation~\ref{eq:bstar}
gives the batch size at which expert offload stops costing throughput, in closed
form, independent of expert size and GPU count, and predicts the measured
crossover to within 0.5\% on OLMoE-1B-7B and to within the resolution of our
measurements on DeepSeek-V2-Lite. MoE-Lightning's HRM is the nearest relative and
selects a micro-batch policy for a system in hand; $B^{*}$ is a threshold a reader
can evaluate for hardware they do not own.

The second new element is the reason behind llama.cpp's design. Unranked static
placement is, on our measurements, the correct choice at serving scale, and we can
now say why: hit rate equals residency, so the ranking that the rest of the
literature computes buys nothing above $B^{*}$, and a ranking computed on the
wrong domain buys less than nothing. What remains is a capacity decision with a
closed form, and a scheduling problem solved by moving whole layers one step ahead
of demand.

These design points are complementary rather than competing, and they partition
cleanly by regime: predict-and-cache below $B^{*}$, contiguous layer streaming
above it. Published figures across the two groups are not directly comparable,
since retention against a resident baseline on a GPU that fits the model measures
a different quantity from speedup against another offloading system on a GPU that
does not. Running these systems on common hardware across the batch range, with
$B^{*}$ predicting where each should win, is the experiment this framework is
built to support.

%  Suggested placement: after the serving-loop section, before
%  "Two smaller results". By that point every measurement this section
%  refers to has already been presented.
% =====================================================================

\section{How this compares to existing expert offloading}

Expert offloading is a well-populated area. Every system below starts from the
same premise we do, that expert weights dominate the footprint and should not
occupy HBM, and then diverges on what to do about it. The divergence is almost
entirely about the batch size each system was designed for, which is the axis our
closed form describes. This section places MEMoE-RT on that axis and states what
its architecture buys as the deployment grows.

\subsection{The design space, ordered by batch size}

\textbf{Predict and cache, at batch one.} The largest group treats expert
selection as predictable and expert reuse as temporally local.
Mixtral-offloading keeps an LRU cache of individual experts on the GPU and
speculatively prefetches one to two layers ahead, on the premise that consecutive
tokens activate overlapping experts (Eliseev and Mazur, arXiv:2312.17238).
MoE-Infinity adds request-level activation tracing, ProMoE trains a predictor,
HOBBIT varies expert precision by importance and compares LRU against
sequence-level LFU and a layer-distance policy, and AdapMoE skips experts judged
unimportant. All of them are built and evaluated at or near single-sequence
decoding on consumer hardware.

\textbf{Move the activation instead of the weight.} Fiddler observes that at
small batch an activation is far smaller than an expert's weights, and therefore
copies activations to the host and executes the expert layer on the CPU rather
than moving weights to the GPU (Kamahori et al., ICLR 2025, arXiv:2402.07033).
Initial placement is popularity-ranked from calibration data. The paper is
explicit that the advantage holds when the batch is sufficiently small, which
makes its design point and ours two sides of the same boundary.

\textbf{Pipeline for throughput.} ZeRO-Inference offloads layer weights to host
or NVMe with prefetch and targets throughput at large batch rather than
single-request latency; its own reported finding is that in that regime full
offload beats partial offload, because the larger batch it enables more than pays
for the transfer (DeepSpeed, 2022). FlexGen schedules tensor placement across
GPU, CPU and disk. MoE-Lightning adds a CPU-GPU-I/O pipeline with paged weights
and a Hierarchical Roofline Model used to select a policy, reporting up to
10.3$\times$ over prior offloading systems for Mixtral-8x7B on a single T4 (Cao
et al., ASPLOS 2025, arXiv:2411.11217). MoE-Gen batches per module.

\textbf{Static placement, in production.} llama.cpp's \texttt{--cpu-moe},
\texttt{--n-cpu-moe} and \texttt{--override-tensor} pin whole routed-expert
tensors to host memory by layer. There is no ranking and no prediction. It is the
simplest arrangement that works and it is what most practitioners actually run.

\subsection{Why the first two groups do not carry to serving scale}

The systems in the first two groups are not underdesigned. They are correctly
designed for $B < B^{*}$, and three of our measurements explain why their
mechanisms stop paying above it.

\emph{Prediction has nothing left to predict.} Coverage saturates at 97.3\% to
100\% by a batch of 1{,}024 tokens (Section~4.2). A speculative prefetcher that
guesses which experts will fire is, at serving batch sizes, guessing that all of
them will, and it is right. The prediction machinery becomes a fixed cost against
a decision that has become trivial.

\emph{Ranking has nothing left to rank.} Fetch hit rate equals resident fraction
exactly at every routing skew we tested at batch 512 and above
(Table~\ref{tab:skewsweep}), which collapses LRU, sampled LFU and static
placement to the same number. An LRU cache above $B^{*}$ is a fixed partition
that pays eviction bookkeeping for nothing. Worse, a placement profiled on one
domain scores 0.383 against 0.500 for random when it meets another
(Section~4.4), so calibration-derived ranking is not merely inert at scale but
can be negative.

\emph{Activation size overtakes weight size.} Fiddler's premise is that
$B \times d_{\text{model}}$ is small against $3 \, d_{\text{model}}
d_{\text{ff}}$. That ratio is linear in $B$, so the premise inverts at a batch
size determined by the same quantities that appear in Equation~\ref{eq:bstar}.
Above the crossover, moving the weights once per layer and computing on the GPU
is the cheaper arrangement, which is what MEMoE-RT does.

\emph{Per-expert granularity fails a concurrency requirement.} Systems that fetch
experts individually issue many small transfers. Section~12.3 gives the standing
requirement: saturating a tier of bandwidth $BW$ at latency $L$ with requests of
size $S$ needs $BW \cdot L / S$ requests in flight, roughly 160 for a 76.8~GB/s
expander at 64-byte granularity. Our transfer-size measurement gives the same
result from the other direction, with time spent waiting rather than moving rising
from 0.16\% for a whole 12~MB expert to 83.18\% at 4~KB. A fetch engine built
around individual experts is committed to a concurrency budget it must then work
to meet.

\subsection{What our architecture does instead}

MEMoE-RT is built from those three findings rather than around them, and each
design decision is a deletion.

Because coverage saturates, there is no predictor. Because ranking is inert, there
is no popularity ordering anywhere in the implementation; the resident set is the
low expert indices. Because latency is negligible at expert granularity and
dominant below it, there is no per-expert fetch path: every non-resident expert
of layer $L+d$ moves as one contiguous transfer, issued $d$ layers ahead on a
dedicated copy stream, into a ring of $d+1$ staging buffers where layer $L$
occupies slot $L \bmod (d+1)$.

That single-descriptor-per-layer structure is what makes the concurrency
requirement disappear rather than be satisfied: one large DMA replaces thousands
of in-flight small reads, and a single expert-sized transfer already reaches 90\%
of the PCIe plateau. It is also what made deep lookahead unnecessary. Our own
simulator, which modelled per-expert on-demand fetches, recommended depth 4 to 8;
the runtime measured depth 1 as sufficient and depth 8 as 41\% more staging memory
for slightly less throughput (Section~7.3). Batching the transfer by layer removed
the need for depth, and we report the correction rather than the original
recommendation.

The second architecture path is the part that generalises furthest. Rather than
reimplement routing per model family, it points each offloaded expert's projection
weights at a view into the staging buffer and drives the tier from pre- and
post-hooks on the MoE block. Because layer $L$ always lands in slot
$L \bmod (d{+}1)$ at a fixed offset, those views are bound once at load time and
remain correct for the process lifetime. The model's forward is untouched. This is
why the same runtime covers both a fused-expert checkpoint and DeepSeek-V2-Lite's
own modelling code, with different top-$k$, shared experts, a dense first layer
and multi-head latent attention, without architecture-specific transfer logic.

\subsection{Scaling behaviour}

The distinction that matters for a deployment is how each design behaves as the
workload grows along four axes.

\begin{table}[htbp]
\centering\small
\caption{Scaling behaviour by axis. The first two groups are built for the
left-hand end of the batch axis and their mechanisms lose value moving right; the
throughput group and MEMoE-RT gain value moving right.}
\label{tab:scaling}
\begin{tabular}{lll}
\toprule
Axis & Predict-and-cache designs & MEMoE-RT \\
\midrule
Batch size & Hit rate $\to$ residency; predictor idles & Cost per token falls as $1/B$ \\
Transfer count & Grows with experts touched & One descriptor per layer, fixed \\
Model size & Bounded by what the GPU holds & Bounded by host memory \\
GPU count & Placement re-profiled per deployment & $B^{*}$ independent of $N$ \\
\bottomrule
\end{tabular}
\end{table}

\textbf{Batch.} A fetched expert serves $Bk/E$ tokens, so expert arithmetic
intensity grows with the batch while the transfer cost is paid once. This is why
the same configuration costs 62\% of throughput at a decode-sized batch and 5.6\%
at a prefill-sized one, and why offloaded decode time per output token is flat at
$325 \pm 2$~ms across an eightfold change in concurrency while resident cost grows
linearly (Section~10.2).

\textbf{Transfer count.} Per-layer batching makes the number of transfers a
function of depth and layer count alone, not of routing, batch size or residency.
The bytes moved scale; the descriptor count does not.

\textbf{Model size.} Because placement is by capacity rather than by profile, the
resident fraction is a free parameter down to zero. DeepSeek-V2-Lite is 31.4~GB
against a 23~GB device, so no resident baseline exists at all, and MEMoE-RT runs
it in 15.13~GB at 7{,}241 tokens per second. Designs that keep a hot set sized
from a profile do not have that degree of freedom in the same way.

\textbf{GPU count.} $B^{*}$ is independent of both expert size and GPU count,
because a larger expert costs proportionally more to move and does proportionally
more work once moved, and because compute and per-node bandwidth scale together
under tensor parallelism. The operating point therefore transfers across
deployment sizes without re-derivation. Pooling is the one place this changes, and
it changes predictably: capacity saving and required batch size both scale with
node count, at one to one.

\subsection{Where the contribution sits}

Several of MEMoE-RT's mechanisms have precedent, and the comparison is more useful
for saying so. Layer-granular prefetch on a separate stream is what
ZeRO-Inference does. Static unranked expert placement is what llama.cpp does. That
full offload beats partial offload in the throughput regime is ZeRO-Inference's
result, not ours. Having an analytical model select a policy is MoE-Lightning's
HRM.

What is new is the boundary itself. Existing systems optimise a point in the
design space and report a speedup against a chosen baseline; Equation~\ref{eq:bstar}
gives the batch size at which expert offload stops costing throughput, in closed
form, independent of expert size and GPU count, and predicts the measured
crossover to within 0.5\% on OLMoE-1B-7B and to within the resolution of our
measurements on DeepSeek-V2-Lite. MoE-Lightning's HRM is the nearest relative and
selects a micro-batch policy for a system in hand; $B^{*}$ is a threshold a reader
can evaluate for hardware they do not own.

The second new element is the reason behind llama.cpp's design. Unranked static
placement is, on our measurements, the correct choice at serving scale, and we can
now say why: hit rate equals residency, so the ranking that the rest of the
literature computes buys nothing above $B^{*}$, and a ranking computed on the
wrong domain buys less than nothing. What remains is a capacity decision with a
closed form, and a scheduling problem solved by moving whole layers one step ahead
of demand.

These design points are complementary rather than competing, and they partition
cleanly by regime: predict-and-cache below $B^{*}$, contiguous layer streaming
above it. Published figures across the two groups are not directly comparable,
since retention against a resident baseline on a GPU that fits the model measures
a different quantity from speedup against another offloading system on a GPU that
does not. Running these systems on common hardware across the batch range, with
$B^{*}$ predicting where each should win, is the experiment this framework is
built to support.

%  Suggested placement: after the serving-loop section, before
%  "Two smaller results". By that point every measurement this section
%  refers to has already been presented.
% =====================================================================

\section{How this compares to existing expert offloading}

Expert offloading is a well-populated area. Every system below starts from the
same premise we do, that expert weights dominate the footprint and should not
occupy HBM, and then diverges on what to do about it. The divergence is almost
entirely about the batch size each system was designed for, which is the axis our
closed form describes. This section places MEMoE-RT on that axis and states what
its architecture buys as the deployment grows.

\subsection{The design space, ordered by batch size}

\textbf{Predict and cache, at batch one.} The largest group treats expert
selection as predictable and expert reuse as temporally local.
Mixtral-offloading keeps an LRU cache of individual experts on the GPU and
speculatively prefetches one to two layers ahead, on the premise that consecutive
tokens activate overlapping experts (Eliseev and Mazur, arXiv:2312.17238).
MoE-Infinity adds request-level activation tracing, ProMoE trains a predictor,
HOBBIT varies expert precision by importance and compares LRU against
sequence-level LFU and a layer-distance policy, and AdapMoE skips experts judged
unimportant. All of them are built and evaluated at or near single-sequence
decoding on consumer hardware.

\textbf{Move the activation instead of the weight.} Fiddler observes that at
small batch an activation is far smaller than an expert's weights, and therefore
copies activations to the host and executes the expert layer on the CPU rather
than moving weights to the GPU (Kamahori et al., ICLR 2025, arXiv:2402.07033).
Initial placement is popularity-ranked from calibration data. The paper is
explicit that the advantage holds when the batch is sufficiently small, which
makes its design point and ours two sides of the same boundary.

\textbf{Pipeline for throughput.} ZeRO-Inference offloads layer weights to host
or NVMe with prefetch and targets throughput at large batch rather than
single-request latency; its own reported finding is that in that regime full
offload beats partial offload, because the larger batch it enables more than pays
for the transfer (DeepSpeed, 2022). FlexGen schedules tensor placement across
GPU, CPU and disk. MoE-Lightning adds a CPU-GPU-I/O pipeline with paged weights
and a Hierarchical Roofline Model used to select a policy, reporting up to
10.3$\times$ over prior offloading systems for Mixtral-8x7B on a single T4 (Cao
et al., ASPLOS 2025, arXiv:2411.11217). MoE-Gen batches per module.

\textbf{Static placement, in production.} llama.cpp's \texttt{--cpu-moe},
\texttt{--n-cpu-moe} and \texttt{--override-tensor} pin whole routed-expert
tensors to host memory by layer. There is no ranking and no prediction. It is the
simplest arrangement that works and it is what most practitioners actually run.

\subsection{Why the first two groups do not carry to serving scale}

The systems in the first two groups are not underdesigned. They are correctly
designed for $B < B^{*}$, and three of our measurements explain why their
mechanisms stop paying above it.

\emph{Prediction has nothing left to predict.} Coverage saturates at 97.3\% to
100\% by a batch of 1{,}024 tokens (Section~4.2). A speculative prefetcher that
guesses which experts will fire is, at serving batch sizes, guessing that all of
them will, and it is right. The prediction machinery becomes a fixed cost against
a decision that has become trivial.

\emph{Ranking has nothing left to rank.} Fetch hit rate equals resident fraction
exactly at every routing skew we tested at batch 512 and above
(Table~\ref{tab:skewsweep}), which collapses LRU, sampled LFU and static
placement to the same number. An LRU cache above $B^{*}$ is a fixed partition
that pays eviction bookkeeping for nothing. Worse, a placement profiled on one
domain scores 0.383 against 0.500 for random when it meets another
(Section~4.4), so calibration-derived ranking is not merely inert at scale but
can be negative.

\emph{Activation size overtakes weight size.} Fiddler's premise is that
$B \times d_{\text{model}}$ is small against $3 \, d_{\text{model}}
d_{\text{ff}}$. That ratio is linear in $B$, so the premise inverts at a batch
size determined by the same quantities that appear in Equation~\ref{eq:bstar}.
Above the crossover, moving the weights once per layer and computing on the GPU
is the cheaper arrangement, which is what MEMoE-RT does.

\emph{Per-expert granularity fails a concurrency requirement.} Systems that fetch
experts individually issue many small transfers. Section~12.3 gives the standing
requirement: saturating a tier of bandwidth $BW$ at latency $L$ with requests of
size $S$ needs $BW \cdot L / S$ requests in flight, roughly 160 for a 76.8~GB/s
expander at 64-byte granularity. Our transfer-size measurement gives the same
result from the other direction, with time spent waiting rather than moving rising
from 0.16\% for a whole 12~MB expert to 83.18\% at 4~KB. A fetch engine built
around individual experts is committed to a concurrency budget it must then work
to meet.

\subsection{What our architecture does instead}

MEMoE-RT is built from those three findings rather than around them, and each
design decision is a deletion.

Because coverage saturates, there is no predictor. Because ranking is inert, there
is no popularity ordering anywhere in the implementation; the resident set is the
low expert indices. Because latency is negligible at expert granularity and
dominant below it, there is no per-expert fetch path: every non-resident expert
of layer $L+d$ moves as one contiguous transfer, issued $d$ layers ahead on a
dedicated copy stream, into a ring of $d+1$ staging buffers where layer $L$
occupies slot $L \bmod (d+1)$.

That single-descriptor-per-layer structure is what makes the concurrency
requirement disappear rather than be satisfied: one large DMA replaces thousands
of in-flight small reads, and a single expert-sized transfer already reaches 90\%
of the PCIe plateau. It is also what made deep lookahead unnecessary. Our own
simulator, which modelled per-expert on-demand fetches, recommended depth 4 to 8;
the runtime measured depth 1 as sufficient and depth 8 as 41\% more staging memory
for slightly less throughput (Section~7.3). Batching the transfer by layer removed
the need for depth, and we report the correction rather than the original
recommendation.

The second architecture path is the part that generalises furthest. Rather than
reimplement routing per model family, it points each offloaded expert's projection
weights at a view into the staging buffer and drives the tier from pre- and
post-hooks on the MoE block. Because layer $L$ always lands in slot
$L \bmod (d{+}1)$ at a fixed offset, those views are bound once at load time and
remain correct for the process lifetime. The model's forward is untouched. This is
why the same runtime covers both a fused-expert checkpoint and DeepSeek-V2-Lite's
own modelling code, with different top-$k$, shared experts, a dense first layer
and multi-head latent attention, without architecture-specific transfer logic.

\subsection{Scaling behaviour}

The distinction that matters for a deployment is how each design behaves as the
workload grows along four axes.

\begin{table}[htbp]
\centering\small
\caption{Scaling behaviour by axis. The first two groups are built for the
left-hand end of the batch axis and their mechanisms lose value moving right; the
throughput group and MEMoE-RT gain value moving right.}
\label{tab:scaling}
\begin{tabular}{lll}
\toprule
Axis & Predict-and-cache designs & MEMoE-RT \\
\midrule
Batch size & Hit rate $\to$ residency; predictor idles & Cost per token falls as $1/B$ \\
Transfer count & Grows with experts touched & One descriptor per layer, fixed \\
Model size & Bounded by what the GPU holds & Bounded by host memory \\
GPU count & Placement re-profiled per deployment & $B^{*}$ independent of $N$ \\
\bottomrule
\end{tabular}
\end{table}

\textbf{Batch.} A fetched expert serves $Bk/E$ tokens, so expert arithmetic
intensity grows with the batch while the transfer cost is paid once. This is why
the same configuration costs 62\% of throughput at a decode-sized batch and 5.6\%
at a prefill-sized one, and why offloaded decode time per output token is flat at
$325 \pm 2$~ms across an eightfold change in concurrency while resident cost grows
linearly (Section~10.2).

\textbf{Transfer count.} Per-layer batching makes the number of transfers a
function of depth and layer count alone, not of routing, batch size or residency.
The bytes moved scale; the descriptor count does not.

\textbf{Model size.} Because placement is by capacity rather than by profile, the
resident fraction is a free parameter down to zero. DeepSeek-V2-Lite is 31.4~GB
against a 23~GB device, so no resident baseline exists at all, and MEMoE-RT runs
it in 15.13~GB at 7{,}241 tokens per second. Designs that keep a hot set sized
from a profile do not have that degree of freedom in the same way.

\textbf{GPU count.} $B^{*}$ is independent of both expert size and GPU count,
because a larger expert costs proportionally more to move and does proportionally
more work once moved, and because compute and per-node bandwidth scale together
under tensor parallelism. The operating point therefore transfers across
deployment sizes without re-derivation. Pooling is the one place this changes, and
it changes predictably: capacity saving and required batch size both scale with
node count, at one to one.

\subsection{Where the contribution sits}

Several of MEMoE-RT's mechanisms have precedent, and the comparison is more useful
for saying so. Layer-granular prefetch on a separate stream is what
ZeRO-Inference does. Static unranked expert placement is what llama.cpp does. That
full offload beats partial offload in the throughput regime is ZeRO-Inference's
result, not ours. Having an analytical model select a policy is MoE-Lightning's
HRM.

What is new is the boundary itself. Existing systems optimise a point in the
design space and report a speedup against a chosen baseline; Equation~\ref{eq:bstar}
gives the batch size at which expert offload stops costing throughput, in closed
form, independent of expert size and GPU count, and predicts the measured
crossover to within 0.5\% on OLMoE-1B-7B and to within the resolution of our
measurements on DeepSeek-V2-Lite. MoE-Lightning's HRM is the nearest relative and
selects a micro-batch policy for a system in hand; $B^{*}$ is a threshold a reader
can evaluate for hardware they do not own.

The second new element is the reason behind llama.cpp's design. Unranked static
placement is, on our measurements, the correct choice at serving scale, and we can
now say why: hit rate equals residency, so the ranking that the rest of the
literature computes buys nothing above $B^{*}$, and a ranking computed on the
wrong domain buys less than nothing. What remains is a capacity decision with a
closed form, and a scheduling problem solved by moving whole layers one step ahead
of demand.

These design points are complementary rather than competing, and they partition
cleanly by regime: predict-and-cache below $B^{*}$, contiguous layer streaming
above it. Published figures across the two groups are not directly comparable,
since retention against a resident baseline on a GPU that fits the model measures
a different quantity from speedup against another offloading system on a GPU that
does not. Running these systems on common hardware across the batch range, with
$B^{*}$ predicting where each should win, is the experiment this framework is
built to support.
